\section*{Abstract}
% -------------------------------------------------------

The linguistic dating of ancient texts from internal morphosyntactic
evidence is central to biblical studies and classical philology, but the
field's methods have been challenged as circular: models are validated
against dates that are themselves partly linguistic judgements, and
validation designs frequently supply a text's expected date to the model
that is then asked to recover it.  The present paper reports a
leakage-free reassessment.  Every out-of-sample result uses an agnostic
prior; feature screening, standardisation and model fitting are performed
inside each cross-validation fold; and significance is established by
permutation of the date labels through the entire pipeline rather than by
parametric assumption.  Two model families are fitted and both are
reported throughout.  Applied to 25 units of the \bhsa{} database spanning
760--167~BCE, the analysis yields a sharply bounded positive result.
Hebrew morphosyntax carries a genuine diachronic signal: the models order
text pairs chronologically with 69.8\% accuracy across 295 pairs
($p = 4 \times 10^{-12}$), and leave-one-out Spearman correlations between
true and predicted dates are $+0.58$ and $+0.49$ in the two families
($p = 0.016$, $0.024$ against a permutation null).  That signal does not,
however, support absolute dating: leave-one-out mean absolute error is
121--156~yr against 137~yr for the trivial predictor that assigns every
text the corpus mean date, and four-way period assignment reaches only
40--48\%.  We show that previously reported holdout accuracies of order
10--30~yr in this literature, including in an earlier version of this
work, arise because each held-out text is dated under a prior centred on
its own scholarly date; for the most tightly anchored texts the linguistic
data contributes under 6\% of posterior precision.  Applied to the
Pentateuchal sources with distribution-free conformal intervals, the
framework supports one substantive conclusion --- the Priestly,
Deuteronomic and JE composites are post-exilic with probability
$\geq 0.92$ under both families --- while declining to place them within a
particular post-exilic period.  The Song of the Sea and Song of Deborah
emerge as the two earliest of nineteen undated units in both families
without having been trained on.  Code and corpora are publicly available
(see Data Availability).
